% HD-R06. GPT-6 Astra Pro, /root; inherited effort; 21 November 2026.
\subsection{HD-R06: complete point-ideal product orbits in dimension six}
\label{sec:point-ideal-sixfold-orbits}

Let $X_i$ be positive-dimensional complex abelian varieties, let
$p_i\in X_i$, and suppose $\sum_i\dim X_i=6$. Start with the actual
perfect complex
\[
                         F=\boxtimes_i I_{p_i}.
\]
Take any finite composition of actual Fourier--Mukai equivalences,
isogeny pullbacks and pushforwards, arbitrary line-bundle twists,
translations, shifts and derived duals. No resulting whole character,
even multiplied by a nonzero rational scalar, has the form
\[
 e^B\bigl(P(\theta)+\lambda t\bigr),\qquad
 B\in\operatorname{NS}(A)_{\mathbf Q},\quad
 \deg P\le6,\quad 0\ne\lambda t\in W_K(A),                 \tag{R06.1}
\]
on an abelian sixfold with compatible polarization $\theta$ and
imaginary-quadratic Weil action of signature $(3,3)$.
The conclusion is a whole-character obstruction for these specified
actual objects and functors; it consequently applies to every coherent
action on such an object.

\paragraph{The exact target contraction ranks.}
We give the dimension-six form of the calculation HD-N14, since its
numerical ranks are different. Over $\mathbf C$ choose the balanced
bases
\[
 \theta=\sum_{i=1}^3(a_i y_i+b_i z_i),\quad
 \tau_+=a_1a_2a_3z_1z_2z_3,\quad
 \tau_-=b_1b_2b_3y_1y_2y_3.
\]
The products are exterior products. A nonzero rational Weil class is
$u\tau_++u'\tau_-$ with $uu'\ne0$. Put
\[
 U_-=\langle\iota_{a_i},y_i\wedge\rangle_{i=1}^3,
 \qquad U_+=\langle\iota_{b_i},z_i\wedge\rangle_{i=1}^3.
\]
These are the two six-dimensional summands of $HH^1(A)$. Assign
weight $+1$ to $a_i,z_i$ and weight $-1$ to $b_i,y_i$.
The operator map from $U_-$ to its Weil output is injective: its
three contractions and three wedges give six distinct monomials of
weight $5$. The corresponding map from $U_+$ has weight $-5$ and
is likewise injective. The polynomial outputs have weights $-1,1$,
so they cannot cancel these maps. Thus $\operatorname{rank}c_v^1=12$.

In $HH^2(A)$ the two exterior-square summands each have dimension
$\binom62=15$. Their Weil outputs have respective weights $4,-4$;
the polynomial outputs have weights $-2,2$. The three groups of
double contractions, contraction--wedge operators, and double wedges
give respectively $3,9,3$ distinct monomials. Thus each exterior
square is injective, and their images are independent of the
weight-zero mixed block.

The mixed block kills the Weil part. For each of the nine pairs
$(i,j)$, its four operators have images, with fixed ordering signs,
\[
 y_i z_jP(\theta),\quad y_i z_jP'(\theta),\quad
 y_i z_jP'(\theta),\quad y_i z_jP''(\theta).
\]
For $0\le k\le4$, the map
$\eta\mapsto\eta\theta^k$ on $\langle y_i z_j\rangle$ is
injective. To recover the coefficient of $y_i z_j$, choose $k$
paired indices other than the $a_i y_i,b_j z_j$ pairs, and read
the monomial with exactly those $k$ holomorphic factors. Its
coefficient is the requested coefficient times a nonzero factorial;
no different pair has that monomial. Different $k$ have different
bidegrees. With $P(z)=\sum_{j=0}^6p_jz^j$ and $m_j=j!p_j$, the
mixed block is consequently nine copies of
\[
 M_P=\begin{pmatrix}
 m_0&m_1&m_2\\m_1&m_2&m_3\\m_2&m_3&m_4\\
 m_3&m_4&m_5\\m_4&m_5&m_6
 \end{pmatrix}.
\]
Multiplication by $e^B$ is the invertible $(1,1)$-shear of the
Hochschild generators proved in HD-N14. It preserves these ranks.
We have therefore proved the exact formulas
\[
 \operatorname{rank}c_v^1=12,\qquad
 \operatorname{rank}c_v^2=30+9\operatorname{rank}M_P
                              \in\{30,39,48,57\}.             \tag{R06.2}
\]
The value $30$ forces every entry of $M_P$ to vanish, and hence
$P=0$.

\paragraph{The complete product list.}
HD-R05 proves the actual factor contraction maps and their
Hochschild--Kunneth multigrading. Apply exactly those maps in total
dimension six. A dimension-one factor gives
$\operatorname{rank}c_F^1=12-\#\{i:\dim X_i=1\}<12$ and is
excluded by (R06.2). For all other factors, put $k_i=\dim X_i$.
The rank formula of HD-R05 becomes
\[
 \operatorname{rank}c_F^2
 =\sum_i r_2(k_i)+4\sum_{i<j}k_i k_j
 =66-\sum_i k_i^2-\#\{i:k_i=2\},                            \tag{R06.3}
\]
where $r_2(2)=1$ and $r_2(k)=k(k-1)$ for $k\ge3$.
The entire list is
\[
\begin{array}{c|r|l}
 (k_i)&\operatorname{rank}c_F^2&\text{remaining argument}\\\hline
 (6)&30&\text{pure Weil character, treated below}\\
 (2,4)&45&\text{excluded by (R06.2)}\\
 (3,3)&48&\text{two-factor spin stabilizer, treated below}\\
 (2,2,2)&51&\text{excluded by (R06.2)}
\end{array}                                                   \tag{R06.4}
\]
One, two and three parts give exactly the displayed partitions;
four parts of size at least two cannot sum to six.
The full actual functor diagrams of HD-R05 preserve both
contraction ranks, so these exclusions hold throughout the stated
orbits, including mixed twists and derived duality.

\paragraph{The product of two threefold point ideals.}
Write $V_i=H^1(X_i,\mathbf R)$, $\dim_{\mathbf R}V_i=6$,
and let $\omega_i$ be its nonzero point volume. The whole source
character is $(1-\omega_1)(1-\omega_2)$. Use the full spin action
on $\bigwedge^*(V_1\oplus V_2)$, with Clifford space
$\mathcal W=V\oplus V^*$ and
$Q(v+\xi,w+\eta)=\xi(w)+\eta(v)$. The three Lie summands act as
\[
        \beta\wedge,\qquad D_T-\tfrac12\operatorname{tr}T,
                      \qquad\iota_\gamma.                    \tag{R06.5}
\]
The degree-two output forces $\beta=0$. The degree-ten output is
$\iota_\gamma(\omega_1\omega_2)$; contraction with this top
volume is an isomorphism on dual bivectors, so $\gamma=0$.
Degree zero forces $\operatorname{tr}T=0$. The off-diagonal
blocks of $T$ acting on $\omega_1,\omega_2$ give respectively
the independent multidegrees $(5,1),(1,5)$, with injective maps
from those blocks. Thus they vanish. The independent volume
coefficients force both diagonal traces to vanish. Conversely
all these trace-zero diagonal blocks kill the character. Its full
real stabilizer is exactly
\[
               \mathfrak{sl}_6(\mathbf R)\oplus
                          \mathfrak{sl}_6(\mathbf R).         \tag{R06.6}
\]

Every target (R06.1) has the explicit faithful
$\mathfrak{su}(3,3)$ action on $H^1(A,\mathbf R)$ preserving
$\theta$ and both Weil determinant lines. Its trace is zero, so
it embeds in (R06.5); the precise twist is
$Z\mapsto e^BD_Ze^{-B}=D_Z-(D_ZB)\wedge$. Therefore a transported
target would give an injection of $\mathfrak{su}(3,3)$ into
(R06.6). Complexification identifies $\mathfrak{su}(3,3)$ with
$\mathfrak{sl}_6(\mathbf C)$: the scalar-extension map and its
Hermitian-adjoint inverse decomposition are exactly those written
in HD-N17, paragraph6. Simplicity of the elementary-matrix Lie
algebra therefore makes any nonzero projection injective.
Both dimensions are $35$, so such a projection is an isomorphism.

For $\mathfrak{sl}_6(\mathbf R)$, its Killing form
$12\operatorname{tr}(XY)$ is positive on the symmetric trace-zero
matrices of dimension $20$ and negative on the skew-symmetric
matrices of dimension $15$. For $\mathfrak{su}(3,3)$,
$12\operatorname{Re}\operatorname{tr}(XY)$ is negative on the
compact block-diagonal subspace of dimension $9+9-1=17$ and
positive on the off-diagonal subspace of dimension $18$.
These orthogonal decompositions give signatures $(20,15)$ and
$(18,17)$. A Lie isomorphism intertwines the adjoint maps and
therefore preserves their trace pairing. The displayed signatures
exclude it. This settles the $(3,3)$ row.

\paragraph{The single sixfold point ideal.}
For this source $v=1-\omega$ on a real space $V$ of dimension
$12$. The degree-two, degree-ten and degree-zero outputs of
(R06.5) force successively $\beta=0,\gamma=0,\operatorname{tr}T=0$.
Every such $T$ kills $v$. Its full real stabilizer is therefore
$\mathfrak{sl}_{12}(\mathbf R)$, of dimension $143$.
By (R06.2), a transported target has $P=0$. We compute its
actual full stabilizer after undoing the invertible $e^B$ shear.

Let $V_{\mathbf C}=U\oplus\overline U$ be the two quadratic-field
eigenspaces on $H^1(A,\mathbf C)$, each of dimension six.
The two Weil volumes have annihilators
\[
 L_+=U\oplus\overline U^*,\qquad
 L_-=\overline U\oplus U^*
          \quad\text{inside }(V\oplus V^*)_{\mathbf C}.       \tag{R06.7}
\]
Indeed wedges from the supporting six-dimensional subspace kill
its top volume, and contractions from the complementary dual kill
it. The two annihilators are transverse maximal isotropic spaces
and are conjugate. The actual map
\[
 \bigwedge^*L_-\longrightarrow\bigwedge^*V_{\mathbf C},
 \ell_1\cdots\ell_j\longmapsto
                  c(\ell_1)\cdots c(\ell_j)\tau_+             \tag{R06.8}
\]
is an isomorphism: writing $L_-=\overline U\oplus U^*$,
its standard exterior monomials independently wedge in the
$\overline U$ basis and contract out the $U$ basis. They give
each exterior basis monomial of $V_{\mathbf C}$ once, with a
nonzero sign. Its vacuum and top lines are precisely the two
Weil-volume lines. In these coordinates the target is
$a\cdot1+b\cdot\mathrm{vol}$ with $ab\ne0$ in dimension twelve.
The degree-two and degree-ten equations again kill the bivector
terms. The opposite scalar actions on the vacuum and top force
trace zero, including for the projective stabilizer. Thus its
complex stabilizer is exactly $\mathfrak{sl}_{12}(\mathbf C)$.

Here is the real form and a literal Hermitian basis. Let $I$ be
the real complex structure on $V$ with eigenspaces
$U,\overline U$, and put $J=\operatorname{diag}(I,-I^*)$ on
$V\oplus V^*$. Then $J^2=-1$, $J$ is skew-adjoint for $Q$,
and its eigenspaces are (R06.7). For the complex vector-space
structure defined by $J$, set
\[
                     h(x,y)=Q(x,y)-iQ(x,Jy).
\]
This is conjugate-linear in the first variable and linear in the
second. Choose a real basis $a_1,b_1,\ldots,a_6,b_6$ with
$b_i=Ia_i$, and its actual dual. Then
$Ja_i=b_i$ and $Ja_i^*=b_i^*$, so
$(a_1,a_1^*,\ldots,a_6,a_6^*)$ is a complex basis of
$V\oplus V^*$. Direct substitution gives the Gram matrix $U^6$.
Its sums and differences, divided by $\sqrt2$, give the diagonal
matrix with six positive and six negative entries.

Preserving the two spaces (R06.7), preserving $Q$, and having
complex trace zero are exactly the real equations for the
stabilizer computed through (R06.8). Thus its real Lie algebra is
$\mathfrak{su}(6,6)$, of dimension $143$. Its Killing signature
is $(72,71)$: the positive off-diagonal blocks have dimension
$2\cdot6^2$, and the negative compact blocks have dimension
$2\cdot6^2-1$. The corresponding signature for
$\mathfrak{sl}_{12}(\mathbf R)$ is $(77,66)$, from symmetric
trace-zero and skew-symmetric matrices. The signatures differ,
so the two actual stabilizers cannot be intertwined. This settles
the last row of (R06.4).

\paragraph{Actual transport and check scope.}
The needed intertwiners are those proved from the actual kernels
in HD-G19 and recorded rationally in HD-N17, paragraph5. In
particular an equivalence has
$Lc(w)L^{-1}=c(fw)$ with its graph isometry $f$; an isogeny
uses $(v,\xi)\mapsto(Tv,(T^{-1})^*\xi)$, and pushforward is
its inverse times the degree. A line bundle uses
$(v,\xi)\mapsto(v-\iota_\xi B,\xi)$. Derived dual conjugates
the two bivector summands by minus one and fixes the middle
summand. These literal maps conjugate the full Lie stabilizers;
the degree-one and degree-two Hochschild versions, including
their duality signs, are HD-R05 (R05.8). Translations and shifts
give the identity and a nonzero scalar on characters. Hence both
the contraction exclusions and the exact Lie obstructions apply
to every allowed finite composition. The primary Fourier sources
are Golyshev--Lunts--Orlov (2001), Proposition4.3.7 and
Lemma4.3.7.1, and Orlov (2002), Theorem2.10, with the exact
original locators given in HD-G19.

This exhausts the specified point-ideal products. It introduces
no assertion about cones, extensions, arbitrary retracts or other
sixfold objects. The exact partition certificate enumerates all
eleven partitions and checks the four displayed ranks; the
universal morphisms and real-form arguments are proved above.
